% Tabla A2-T2. Generada automáticamente; no editar a mano.
% Requiere \usepackage{booktabs} (y graphicx si la tabla es ancha).
\begin{table}[htbp]
    \centering
    \caption{Análisis II: comparación de los modelos de la masa planetaria en la muestra de ajuste. Diferencias respecto al modelo base M1.}
    \label{tab:A2-T2}
    \small
    \resizebox{\textwidth}{!}{%
    \begin{tabular}{llrrrrrrrrrrr}
    \toprule
    Modelo & Predictores & $N_{\rm pl}$ & $R^2$ & $R^2_{\rm aj}$ & AIC & BIC & RMSE & MAE & $\Delta R^2$ & $\Delta$AIC & $\Delta$RMSE & $\Delta$MAE \\
    \midrule
    M0 & s(log\_Rp) & 535 & 0.482 & 0.478 & 126.9 & 148.3 & 0.270 & 0.194 & -0.0184 & +17.4 & +0.0048 & +0.0031 \\
    M1 & $s($$\log R_p$$)$ + $\log F$ & 535 & 0.501 & 0.496 & 109.5 & 135.2 & 0.265 & 0.191 & +0.0000 & +0.0 & +0.0000 & +0.0000 \\
    M2 & $s($$\log R_p$$)$ + $\log F$ + $M_\star$ + {[Fe/H]} & 535 & 0.506 & 0.499 & 108.3 & 142.5 & 0.264 & 0.190 & +0.0049 & -1.2 & -0.0013 & -0.0007 \\
    \bottomrule
    \end{tabular}
    }
    \par\smallskip
    {\footnotesize\textit{Nota:} 535 planetas alrededor de 365 estrellas. s() denota un spline cúbico. RMSE y MAE en dex.}
\end{table}
